\section{Secure Bindings and Comparative Intelligence Learning}
\label{sec:aion-flow-secure-bindings-comparison}

This section records completion of AION Flow task \texttt{AFLOW-0211} and
Stage~9.  Together they let a customer connect replaceable intelligence without
placing secrets in a workflow, then compare alternative models, harnesses or
compute routes without surrendering the accumulated AION brain.

\subsection{Visual Vault Credential Binding}

Every Intelligence Stack field requiring a credential is rendered as a Visual
Vault selector.  The selector displays only non-secret metadata: a human label,
provider, stable reference, status and revision.  A new credential may be
entered from the Node Editor, but it is sent directly to the local mother-brain
vault.  The vault encrypts it and returns a stable opaque reference of the form
\texttt{vault://model/binding-name}.  Only this reference is stored in the
canvas graph.

Raw keys, passwords and tokens are therefore excluded from canvas state,
autosave, workflow history, logs, previews, exports and execution receipts.  A
create operation requires an explicit vault-intent marker.  Public list
operations cannot return secret material.  Bindings can subsequently be
rotated or revoked without editing every workflow that refers to them; a
revoked or unavailable binding fails closed.

\subsection{Controlled Stack Comparison}

The existing Workflow Canvas now exposes a \emph{Compare stacks} centre.  A
user chooses one model, harness or compute node and supplies its challenger
value.  AION duplicates the complete graph and changes only the selected
declared field.  The original graph remains the champion.  Both graph versions
are hashed and bound to the same versioned evidence pack, preventing an
apparently superior route from being tested on easier evidence.

The comparison records six normalized dimensions:
\begin{itemize}
  \item quality and correctness;
  \item evidence coverage;
  \item latency;
  \item monetary cost; and
  \item energy consumption.
\end{itemize}
Public benchmark measurements and customer-private outcome measurements use
separate scopes.  Private measures remain in the customer-owned evaluation
ledger and are not promoted into a public benchmark.

\subsection{Simulation, Evidence and Learning Boundary}

The two immutable variants can be run through a paired dry comparison.  Both
receive the identical evidence-pack object and hash.  This operation performs
no external writes and cannot alter route preferences.  It exists to expose
configuration, policy, evidence and cost differences before live evaluation.

A measurement may influence AION only when it carries either a verified
outcome receipt or an explicit human correction.  Unverified model scores and
ordinary simulations are retained as observations but excluded from the
learning sample.  This prevents a model from declaring itself successful and
then teaching the router to select itself.

\subsection{Champion, Challenger and Reversible Promotion}

Challenger traffic is bounded to at most one half of eligible work and defaults
to ten per cent.  Promotion requires the configured minimum sample count,
improved quality, no regression in correctness or evidence coverage, and no
active drift warning.  Drift checks currently identify material correctness or
evidence regression and changes in provider-behaviour fingerprints.

Every promotion stores the actor, reason, comparison hash and decision hash.
The interface explains the accepted learning evidence, ignored unverified
observations, current recommendation, active route, detected drift and routing
history.  The owner may return immediately to the original champion; rollback
does not delete the challenger or the evidence that justified either decision.

\subsection{Sovereign Result}

Models, harnesses, endpoints and compute providers remain replaceable tools.
Identity, memory, permissions, Business Map knowledge, evidence policy,
verified outcomes and learned route preferences remain with the customer's
AION brain.  A provider can therefore be upgraded, compared, demoted or removed
without resetting what the organisation has learned.
